\paperhead
{Optimized Topic Modeling of Nepal-Affiliated AI Research:\\
Comparing Search Strategies and Tracking Themes, 2015--2025}
{%
Mapping a regional research literature requires models whose topics are both measurable and
meaningful. We analyse 1,074 scholarly records retrieved as Nepal-affiliated AI research from
2015--2025, comparing grid-selected latent Dirichlet allocation (LDA), genetic-algorithm-optimised
LDA, a budget-matched random-search control, and BERTopic. GA-LDA improves coherence over the
standard baseline ($C_v$: 0.4483 to 0.4940), but random search finds a better solution in the
single matched-budget experiment (fitness 0.5282 versus 0.5196). BERTopic achieves the highest
$C_v$ (0.5107) and three-seed topic stability (0.8418), while leaving 17.1\% of documents
unassigned. Three blinded LLM judgments favour BERTopic and rate GA-LDA below standard LDA,
showing that improved coherence need not coincide with improved automated interpretability.
An additional label audit identifies several GA topics as generic rather than domain-specific.
Fixed-model temporal analysis shows increasing education-related topic share after 2020 under
both representations. However, unequal window sizes, differing assignment coverage and a
residual non-AI record limit interpretation, particularly for the eight-record earliest window.
The results support reporting multiple quality measures, explicit coverage and qualitative
checks rather than selecting a model from coherence alone.%
}
{topic modelling, latent Dirichlet allocation, genetic algorithms, BERTopic, Nepal, scientometrics}

\section{Introduction}
\label{sec:t1-intro}

A growing publication list does not by itself explain what a research community studies.
Researchers choosing a project, institutions planning research support and readers entering a
field need to distinguish recurring applications from emerging themes. Reading individual
papers remains essential, but an overview becomes difficult when relevant work is scattered
across engineering, health, education, agriculture and other disciplinary venues. Topic
modelling offers a complementary way to organise this literature: it identifies patterns in
document text that can be inspected, labelled and compared across publication periods.

This study examines research with at least one Nepal-based institutional affiliation. That
definition concerns the affiliation recorded on a publication, rather than an author's
nationality or ethnicity. It includes internationally co-authored work when a Nepal affiliation
is present and excludes diaspora-only work without one. The distinction matters because
``research in Nepal'' can otherwise refer to authorship, study location, language or national
identity. Our evidence supports a view of indexed Nepal-affiliated literature, not a complete
census of all AI activity in the country.

Latent Dirichlet allocation (LDA)~\cite{blei2003} represents documents as mixtures of topics and
topics as distributions over words. Its use in scientific literature analysis~\cite{griffiths2004}
provides an interpretable starting point for this application. Large-scale AI mapping already
includes temporal and country-level comparisons~\cite{yu2023}; the question here is how useful
different topic models are within a smaller regional corpus with heterogeneous applications and
uneven temporal coverage. We investigate this particular setting without claiming that Nepal
has never been studied or that regional analysis is itself novel.

Three methodological difficulties motivate the comparison. First, topic count and prior settings
can change the granularity of an LDA summary: a broad topic may score consistently while failing
to name a recognisable research area. Second, an embedding-based model may separate specific
domains while declining to assign atypical documents, complicating comparisons with a model
that represents every record. Third, publication growth and database coverage can be mistaken
for thematic change unless topic prevalence is normalised and sample sizes remain visible.
These are practical model-selection questions, not simply a contest to maximise one score.

We therefore compare a conventional LDA baseline, an evolutionary hyperparameter search,
a random-search control and BERTopic on the same frozen corpus. The study contributes a
documented regional dataset, an explicit test of whether the evolutionary mechanism adds value
beyond search effort, and a temporal interpretation qualified by assignment coverage and topic
quality. The genetic algorithm is an application of established methods rather than a new
optimisation algorithm. The paper first positions the comparison in related work, then describes
the corpus and methods, before examining quantitative results alongside the words and documents
that give those results meaning.

\section{Related Work}
\label{sec:t1-related}

\subsection{Mapping scientific literature}
Griffiths and Steyvers~\cite{griffiths2004} demonstrate how probabilistic topics can organise
scientific publications. Yu and Xiang~\cite{yu2023} apply LDA to 177,204 AI documents from
1990--2021, extracting 40 topics and aggregating them by year, publication source and country or
region. Their work establishes that topic prevalence can provide a structured account of AI
research. Our contribution differs in scale and emphasis: a Nepal-affiliated corpus is analysed
through competing representations and explicit evaluation of their limitations. A global model's
topic vocabulary cannot be assumed to offer the most useful resolution for this narrower collection.

\subsection{Search and embedding-based topic models}
Agrawal, Fu and Menzies~\cite{agrawal2018} investigate instability in topic modelling and use
search-based optimisation to improve model configurations. This motivates searching LDA's
hyperparameters instead of fixing them arbitrarily. However, improvement over a small grid does
not isolate the contribution of an evolutionary algorithm. Random search is an established
hyperparameter-selection baseline~\cite{bergstra2012}; matching its space, objective and number
of unique evaluations to the GA makes the comparison more informative.

BERTopic~\cite{grootendorst2022} combines document embeddings with clustering and class-based
TF-IDF topic representations. Its components include Sentence-BERT~\cite{reimers2019},
UMAP~\cite{mcinnes2018} and HDBSCAN~\cite{campello2013}. Capilitan et al.~\cite{agri2026}
illustrate its application to a large agricultural AI literature, showing the relevance of
embedding-based thematic mapping to an applied domain. Its usefulness on our collection is an
empirical question: density clustering can discover specific groups but may leave some records
outside those groups. We use hard assignments for temporal reporting; this is a choice in our
implementation, not a claim that BERTopic cannot produce probability estimates.

\subsection{Evaluation beyond a single score}
Chang et al.~\cite{chang2009} show why statistical fit and human interpretation should not be
treated as equivalent. Word co-occurrence approaches~\cite{newman2010,mimno2011} and the
$C_v$ framework of R\"oder et al.~\cite{roder2015} provide measurable aspects of topic quality,
while topic diversity~\cite{dieng2020} addresses repeated vocabulary across topics. None alone
establishes that a topic names a useful research domain. We therefore combine coherence,
diversity, stability, document coverage and blinded automated judgments. The latter are explicitly
provisional: this experiment does not replicate a human interpretation study. The combined
evaluation asks which compromises each model makes and whether those compromises suit a
regional literature review.

\section{Problem and Dataset}
\label{sec:t1-data}

\subsection{Research questions}
\textbf{RQ1:} Which themes are identifiable in the retrieved Nepal-affiliated literature?
\textbf{RQ2:} How does their prevalence change between publication windows?
\textbf{RQ3:} Does GA-LDA improve on grid-selected LDA and on budget-matched random search?
\textbf{RQ4:} How do these configurations compare with BERTopic in quality and coverage?

\subsection{Retrieval and inclusion}
Records were retrieved from OpenAlex~\cite{priem2022} using a Nepal-affiliation country filter,
AI-related topic subfields and a supplementary keyword pass for applied work indexed elsewhere.
The inclusion rule required a Nepal-based institutional affiliation and AI-related vocabulary in
the title or abstract. The analysis uses concatenated titles and abstracts rather than full texts.
The filtering sequence is reported in Table~\ref{tab:funnel}.

\begin{table}[t]
\caption{Frozen corpus inclusion sequence.}
\label{tab:funnel}\centering\small
\begin{tabular}{@{}lr@{}}\toprule
Stage & Records remaining\\\midrule
OpenAlex retrieval & 1,693\\
Publication year through 2025 & 1,228\\
Abstract available & 1,102\\
At least 30 title/abstract words & 1,088\\
Exclude computer vision syndrome & 1,074\\\bottomrule
\end{tabular}
\end{table}

OpenAlex's structured affiliation metadata made it suitable for the geographical inclusion rule.
This departs from the proposal's planned Semantic Scholar-led, multi-source collection: NepJOL,
IEEE and Springer were not separately harvested into this analysis. Indexing through OpenAlex
may nevertheless include records from those publishers or outlets. The absence of direct
supplementation limits coverage; it is not proof that every local publication is absent, nor that
supplementary harvesting is technically impossible.

The corpus contains 560 distinct nonempty venue names and a median title/abstract length of
212 words. Venue metadata classifies 1,004 records as published, 36 as preprint-server records
and 34 as lacking a venue. Thus, 93.5\% have publication-venue evidence, but the dataset is not a
verified wholly peer-reviewed collection as originally proposed. Recorded affiliation strings
are available for 1,072 records. An earlier 30-record affiliation check reported no incorrect
Nepal affiliations; that limited check does not establish complete AI relevance or coverage.

\begin{figure}[t]\centering
\includegraphics[width=\columnwidth]{figCorpusByYear.png}
\caption{Publication counts in the frozen corpus. The last window contains 674 of 1,074 records;
growth reflects both research activity and retrieval coverage.}
\label{fig:years}\end{figure}

\subsection{Quality checks and temporal coverage}
The four windows contain 8, 93, 299 and 674 records respectively (Fig.~\ref{fig:years}). These
counts make the earliest window unsuitable for a separately fitted topic model. Fourteen
``computer vision syndrome'' papers were removed because their use of that phrase concerned
screen-related eye conditions rather than computer vision methods. This illustrates why keyword
matching requires inspection of context.

A subsequent document check found another ambiguity: a 2016 medical-education paper discusses
students' ``deep learning'' approaches, not machine learning. It remains in the archived training
corpus, so the comparison is reported as a study of the frozen retrieval rather than a perfectly
adjudicated AI dataset. We record its identifier and quantify its influence on the earliest window
in Section~\ref{sec:temporal}. Policy and ethics papers about AI are also retained. These
qualifications prevent apparent thematic precision from concealing uncertainty in inclusion.

Affiliation validity and subject relevance are separate questions. An institution can be correctly
identified while the record is included for the wrong sense of a technical phrase. A useful
adjudication process therefore checks the abstract's methods as well as title keywords, records
the reason for each disputed inclusion, and preserves identifiers for reproducible exclusions.
Here the residual error is disclosed rather than removed only from the model with an unfavorable
score. Revising the training corpus would require rerunning the entire comparison, since even one
document can change vocabulary frequencies, cluster membership or selected hyperparameters.

\section{Methods}
\label{sec:t1-methods}

Figure~\ref{fig:pipeline} summarises the shared corpus and model-specific representations.
\begin{figure*}[t]\centering
\includegraphics[height=.60\textheight,trim=95 60 95 60,clip,keepaspectratio]{figPipeline_corrected.png}
\caption{End-to-end analysis. All models use the same records, while representations differ to suit
the model. Label assessment and temporal interpretation follow, rather than define, training.}
\label{fig:pipeline}\end{figure*}

\subsection{Text preparation and feature extraction}
Cleaning normalises Unicode, removes leading abstract markers, URLs and DOI strings, and
consolidates parenthesised acronym glosses. LDA's branch lowercases and tokenises the text,
removes stop words, applies noun/verb lemmatisation and detects bigram/trigram phrases. Technical
acronyms such as AI, ML, NLP and CNN are protected from the minimum-length filter and lemmatiser.
The dictionary retains terms appearing in at least five documents and at most half the corpus,
yielding 3,018 terms and no empty documents. LDA receives integer bag-of-words counts; the
pipeline also exports TF-IDF features, but these are not used to train the probabilistic model.

The domain stop-word list removes publication boilerplate and frequent generic research-process
terms. Its development used measured document frequency alongside semantic judgment; frequency
alone cannot distinguish a generic term from an important subject. For example, ``level'' and
``value'' were retained because they can carry hydrological or statistical meaning. Terms such as
``data'', ``accuracy'' and ``performance'' also remain. Their persistence in some topics is therefore
a limitation of the representation, not evidence that cleaning removed them successfully.

BERTopic receives cleaned natural sentences for embedding. Removing function words and word
order before sentence encoding would discard information the encoder expects. The topic-word
representation uses the configured stop words, but tokenisation and phrase vocabularies still
differ between branches. Shared documents and a shared stop list do not make preprocessing
identical; this difference is considered when interpreting coherence.

\subsection{LDA baseline and genetic search}
LDA is fitted with Gensim~\cite{rehurek2010} using variational inference~\cite{hoffman2010}.
The baseline selects $K$ from $\{5,10,\ldots,40\}$ by maximum $C_v$, with symmetric document
priors and the library's default topic-word prior. The selected value is $K=10$
(Fig.~\ref{fig:coherence-k}); it is a searched baseline, not an arbitrary topic count.

\begin{figure}[t]\centering
\includegraphics[width=\columnwidth]{figCoherenceVsK.png}
\caption{Baseline topic-count sweep. Selection uses $C_v$ only; the displayed diversity is D@10
and is a diagnostic rather than a selection criterion.}
\label{fig:coherence-k}\end{figure}

The genetic algorithm~\cite{holland1975} searches $(K,\alpha,\eta)$, where Gensim's $\eta$ is
the topic-word Dirichlet prior often denoted $\beta$. Topic count ranges from 5 to 40, and each
scalar prior ranges from 0.01 to 1.0. Initial priors are sampled log-uniformly to cover sparse
and dense settings. Chromosomes use an integer topic count and real-valued priors. Tournament
selection, blend crossover, Gaussian mutation and two-member elitism generate subsequent
populations. Mutation is scaled to the gene range; bounds and valid integer topic counts are
enforced after variation.

Fitness scalarises three objectives:
\begin{equation}
 \begin{aligned}
 F&=0.6\widehat{C_v}+0.2D_{25}+0.2S_{10},\\
 \widehat{C_v}&=\operatorname{clip}\!\left(\frac{C_v-0.30}{0.40},0,1\right).
 \end{aligned}
\end{equation}
Here $D_{25}$ is the fraction of unique words across top-25 lists, and $S_{10}$ is matched
top-10-word stability. The weights express a preference for coherence; they are not learned and
no Pareto frontier is estimated. Alternative weights could select different configurations.
The historical GA evaluates 20 individuals over 15 generations; caching duplicate genomes
reduces 300 candidate requests to 183 unique evaluations, each requiring three LDA fits.

\subsection{Random search and BERTopic}
Random search uses the same hyperparameter bounds, log-uniform initial sampling rule, fitness,
LDA settings and stability seeds. Its budget is 183 unique configurations, matched to the GA's
actual evaluations. This controls for searching a larger space than the eight-point baseline.
The strategies differ in how subsequent candidates are proposed; matching evaluations does not
imply identical wall-clock overhead or equal performance over repeated search runs.

BERTopic uses \texttt{all-MiniLM-L6-v2} embeddings, UMAP with 15 neighbours and five output
dimensions, and HDBSCAN with minimum cluster size 15 and excess-of-mass cluster selection.
Its 14 non-outlier clusters provide the topics used here. The outlier label is retained in
coverage calculations instead of being reassigned to improve apparent completeness.

\subsection{Stability and temporal summaries}
Across each pair of fitted seeds, topic word sets are matched one-to-one by maximising total
Jaccard overlap with Hungarian assignment. The score averages matched overlap across topics and
seed pairs, penalising unequal topic counts by their size ratio. This avoids treating arbitrary
topic indices as meaningful. It measures reproducibility of top words, not agreement of document
assignments or complete topic distributions.

One global model defines each topic vocabulary across all years. LDA's temporal share is the mean
document-topic mass within a window. BERTopic's share is the fraction of assigned documents in a
cluster; zero-mass outlier rows are excluded. Its denominator therefore differs from LDA's. We
show assigned and total counts for every window and exclude 2015--2017 from trend fitting.
This replaces the proposal's per-window refitting: prevalence changes are measured under fixed
topics, without claiming to identify topic births or changes in topic meaning.

\subsection{Automated thematic labelling}
All 39 topics have archived LLM-assisted draft labels, followed by an automated label audit.
The drafts were produced in an assistant session, not through provider API calls. This fulfils
the proposal's automated-labelling step; it does not establish human validation.

A Streamlit app exposes the process through a wrapper for OpenAI,
Claude and Gemini. It previews top words and up to four representative titles, generates
structured labels with rationale and evidence identifiers, and records prompts, source hashes,
requested/resolved models, responses and usage. Random-search topics use words only because
their archived run lacks document-topic assignments. Failed batches can resume; named human
reviews and partial exports remain explicit. Provider adapters were tested with simulated
responses; no new live API run was made because credentials were unavailable. The reported
results retain their archived labels and judgments.

\section{Experimental Setup}
\label{sec:setup}

The archived search seed is 42, with seeds 42--44 used for stability. LDA uses ten passes,
100 iterations and batches of 500 documents. GA settings are population 20, 15 generations,
tournament size three, crossover probability 0.7, mutation probability 0.05 and mutation scale
0.15 times the gene range. Coherence is evaluated in a single process. The comparison reports
$C_v$, $c_{\mathrm{npmi}}$ and $u_{\mathrm{mass}}$ on top-10 words, and diversity/Jaccard on
top-10 lists. \textbf{The reported D@10 differs from D@25 inside search fitness}; both cutoffs
are explicit so the table is not mistakenly read as the fitness decomposition.

Perplexity is evaluated on the training corpus and is an in-sample variational diagnostic.
Metrics not used in fitness provide additional views, but are not independent held-out tests.
The original three-seed stability also reuses selection seeds. A separate fixed-configuration
check on seeds 100--109 is reported as a sensitivity analysis, not as repeated optimisation or
out-of-sample accuracy. No chronological prediction task is claimed.

For BERTopic coherence, words missing from the shared LDA dictionary are dropped; 93 of 140
topic words remain. The resulting scores therefore compare different effective word-list
lengths, an important limitation of numerical ranking. A topic with fewer than two retained
words would be omitted by the scorer. Embedding-based topic quality should not be inferred
from the filtered score alone.

The interpretability instrument presents shuffled top-10 lists and a 1--5 rubric to three
blinded LLM judgments. Model names, candidate labels and other judgments are withheld. A
separate automated reviewer assesses all 39 candidate labels as accept, revise or reject.
The records identify the judging model as \texttt{claude-opus-5}; separate contexts are recorded,
but sampling settings are not fully preserved. Thus the aggregation is reproducible from saved
ratings, whereas an identical fresh judgment cannot be guaranteed. Ordinal Krippendorff's
$\alpha$~\cite{krippendorff2004} measures agreement. Shared model biases remain possible, and
these ratings are not a substitute for human validation.

To assess sampling sensitivity of the education trend, we resample documents within the
2018--2020 and 2024--2025 windows 2,000 times and take percentile intervals for the difference
in mean share. These intervals condition on the fitted model and, for BERTopic, assigned
documents. They exclude uncertainty from retrieval, topic fitting, labelling and outlier selection.

\section{Results}
\label{sec:results}

\subsection{Model comparison and search behaviour}
\begin{table*}[t]
\caption{Archived model comparison. Arrows indicate preferred metric direction, not universal
model utility. Stability uses the original three seeds; perplexity is in-sample.}
\label{tab:comparison}\centering\small
\begin{tabular}{@{}lrrrr@{}}\toprule
Metric & Standard LDA & GA-LDA & Random-search LDA & BERTopic\\\midrule
Topics & 10 & 7 & 8 & 14\\
$C_v\uparrow$ & .4483 & .4940 & .5030 & \textbf{.5107}\\
$c_{\mathrm{npmi}}\uparrow$ & -.0230 & .0112 & \textbf{.0245} & -.0300\\
$u_{\mathrm{mass}}\uparrow$ & -3.3586 & \textbf{-2.1796} & -2.5477 & -3.4444\\
Diversity D@10 $\uparrow$ & .7900 & \textbf{.8000} & \textbf{.8000} & .7286\\
Within-model Jaccard J@10 $\downarrow$ & \textbf{.0359} & .0471 & .0536 & .0471\\
Stability S@10 $\uparrow$ & .2526 & .3375 & .3382 & \textbf{.8418}\\
In-sample perplexity $\downarrow$ & 172.1 & 167.8 & \textbf{165.9} & n/a\\
Unassigned documents & 0\% & 0\% & 0\% & 17.1\%\\
Mean LLM rating (1--5) & 2.57 & 2.05 & 2.88 & \textbf{4.10}\\
Labels accepted unchanged & 5/10 & 1/7 & 5/8 & 9/14\\\bottomrule
\end{tabular}\end{table*}

For RQ3, GA-LDA improves $C_v$ by 10.2\% relative to the grid baseline and also improves the
other coherence measures, reported diversity, original-seed stability and in-sample perplexity
(Table~\ref{tab:comparison}). It does not improve every metric: standard LDA has less pairwise
topic overlap and a higher automated interpretability score. Coverage is unchanged.

Random search finds the better scalarised fitness, 0.5282 versus 0.5196, with $K=8$ rather than
the GA's $K=7$. Figure~\ref{fig:search} compares incumbent fitness at the same cumulative number
of unique evaluations. The GA stops improving at generation eight, whereas the random run
eventually finds one candidate exceeding the GA result. This supports a conclusion about the
observed runs only. It does not establish that random search generally wins, that GA reduces
variance, or that evolutionary search contributes nothing in other settings.

\begin{figure}[t]\centering
\includegraphics[width=\columnwidth]{figSearchBudget.pdf}
\caption{Best-so-far search fitness against actual unique evaluations, reconstructed from the
saved candidate histories. Each curve represents one search run; no uncertainty band is implied.}
\label{fig:search}\end{figure}

For RQ4, BERTopic leads in $C_v$ and original-seed stability but has lower diversity and poorer
$c_{\mathrm{npmi}}$ and $u_{\mathrm{mass}}$. It leaves 184 documents unassigned. Its top words
are comparatively reproducible across the three seeds, but filtered vocabulary and different
representations qualify direct metric comparisons. No single row resolves the practical choice
between complete document representation and more specific clusters.

On ten additional seeds (100--109), matched-word stability for the fixed standard LDA,
GA-LDA and random-search configurations is 0.2718, 0.3525, 0.3581,
respectively. The same ordering persists, but absolute LDA stability remains modest. These
45 pairwise comparisons per configuration share fitted runs and are not 45 independent
experiments; no significance test between search strategies is inferred from them.


\subsection{Interpreting the topics}
The panel's ordinal agreement is $\alpha=0.9787$, with 93.2\% pairwise exact agreement and all
rating pairs within one point. The ordering is BERTopic, random search, standard LDA, then
GA-LDA (Fig.~\ref{fig:quality}). This is \emph{not} a complete reversal of $C_v$: the first two
models retain their ranks, while the two remaining LDA configurations exchange positions.
The narrower finding is useful: GA-LDA's coherence gain over standard LDA does not translate
into a higher rating under this automated rubric.

\begin{figure}[t]\centering
\includegraphics[width=\columnwidth]{figQualityTradeoff.pdf}
\caption{Model-level coherence and automated interpretation scores. Four descriptive points do
not establish a correlation or causal effect of optimisation.}
\label{fig:quality}\end{figure}

\begin{figure*}[t]\centering
\includegraphics[width=.95\textwidth]{figTopicExamples.pdf}
\caption{Actual top-10 lists illustrate what aggregate scores can hide: a generic evaluation
topic, an education-related topic and a domain-specific drug-discovery cluster.}
\label{fig:examples}\end{figure*}

Figure~\ref{fig:examples} shows why. GA topic 5 is dominated by evaluation and prediction
vocabulary and cannot confidently be called a crop or hazard theme. Other clusters name more
specific domains. The label audit accepts 20 labels, revises 12 and rejects seven
(Fig.~\ref{fig:labels}). Only one of seven GA labels is accepted unchanged, compared with nine
of fourteen BERTopic labels. This separate task corroborates a difficulty in assigning useful
labels, although its judgments can share biases with the rating panel. Broad topics and low
ratings are associated in this comparison; a controlled topic-count ablation would be needed
to attribute the difference to $K$ alone.

\begin{figure}[t]\centering
\includegraphics[width=\columnwidth]{figLabelAudit.pdf}
\caption{Automated review outcomes, with counts shown rather than treating models with different
numbers of topics as equally sized samples. Human label validation remains outstanding.}
\label{fig:labels}\end{figure}

\subsection{Themes and temporal prevalence}
\label{sec:temporal}
The audited labels identify applied domains including medical imaging, drug discovery,
cybersecurity, construction materials, water/land modelling and Nepali text processing. These
provide an answer to RQ1 at the resolution supported by the retrieved corpus. Mixed GA topics
remain labelled as mixed; they are not recast as specific applications to make a cleaner narrative.

\begin{figure*}[t]\centering
\includegraphics[width=.98\textwidth]{figTemporalGA.pdf}
\caption{GA-LDA topic mass by window, using audited labels and all seven topics. Values are
percentages over all documents. The dashed boundary separates the eight-record descriptive window.}
\label{fig:temporal-ga}\end{figure*}
\begin{figure*}[t]\centering
\includegraphics[width=.98\textwidth]{figTemporalBT.pdf}
\caption{BERTopic cluster prevalence conditional on assignment. Assigned counts are 5, 72, 249
and 564; total counts are 8, 93, 299 and 674. These denominators differ from Fig.~\ref{fig:temporal-ga}.}
\label{fig:temporal-bt}\end{figure*}

For RQ2, GA's education-related share rises from 2.1\% in 2018--2020 to 15.2\% in 2024--2025,
an increase of approximately 13.1 percentage points. BERTopic's education cluster rises from
1.4\% to 19.5\% of assigned records, approximately 18.1 points. Conditional bootstrap intervals
are 10.7--15.4 and 13.7--22.2 points respectively. Both fitted representations therefore indicate
an education-related increase across these endpoints, although the values are not directly
interchangeable and the intervals do not address shared coverage bias.

\begin{figure*}[t]\centering
\includegraphics[width=.95\textwidth]{figCoverageSensitivity.pdf}
\caption{Assignment coverage and a denominator sensitivity check for BERTopic's education
cluster. The grey region marks the eight-record early window. Counting outliers in the total
reduces the estimated level, while the later-window increase persists.}
\label{fig:coverage}\end{figure*}

Figure~\ref{fig:coverage} makes this denominator effect explicit. BERTopic assigns 62.5\% of
the earliest records, 77.4\% in 2018--2020, 83.3\% in 2021--2023 and 83.7\% in 2024--2025.
Its changing coverage could therefore influence a temporal comparison conditioned on assignment.
When the education cluster is instead expressed as a fraction of \emph{all} retrieved documents,
including outliers in the denominator, its share rises from approximately 1.1\% to 16.3\%
between the later endpoints. The increase persists under this alternative summary, although
outliers still have unknown thematic membership. Counting them as not assigned to education
does not establish that none is about education; it measures the observed cluster's corpus share.

This check separates two possible explanations for the displayed rise: a greater fraction of
papers being assigned somewhere, and a larger observed education cluster. It strengthens the
descriptive finding without implying that the cluster is a complete or unbiased inventory of
educational AI. Reading its representative documents and adjudicating unassigned records would
be needed to assess that stronger claim. The same distinction should accompany any proposed
comparison of thematic prevalence across model families.

The earliest GA education share requires particular caution. The non-AI medical-education
record receives 0.721 education-topic mass. Excluding that record from reporting alone reduces
the earliest share from about 10.8\% to 2.1\%, without retraining. This is a sensitivity check,
not a corrected model fit. It demonstrates that a modern label applied retrospectively to a
fixed topic cannot establish early generative-AI activity. The entire earliest window remains
descriptive rather than evidence for a decade-long trajectory.

Declining relative shares for some image/text and mixed forecasting topics do not imply falling
publication counts: the corpus grows sharply. Likewise, a broad education cluster cannot be
equated wholly with generative AI. The defensible interpretation concerns topic prevalence in
the retrieved literature, not a causal account of changes in national research priorities.

\section{Social, Ethical, Legal and Professional Considerations}
\label{sec:selp}
The study concerns an identifiable research community. Affiliation is used for inclusion, not
to infer nationality or rank individual researchers. Findings are reported at topic level;
neither publication counts nor topic scores are measures of a person's research quality.

OpenAlex supplies openly available scholarly metadata. The collection uses metadata and
abstracts rather than bypassing paywalls, and the scraper records source identifiers with
rate-limited requests. Metadata availability should not be read as blanket permission to
redistribute every underlying text; provenance and applicable source terms remain relevant.
Direct supplementation of locally published and Nepali-language scholarship was not completed,
which may underrepresent precisely the work of local interest.

LLMs assisted topic labelling, automated assessment and manuscript preparation. Their outputs
are fallible: generated labels were revised or rejected, and the current report distinguishes
automated judgments from human review. The authors remain responsible for the submission,
its sources and its interpretation. Claims of perfectly independent judges or guaranteed fresh
reproduction of a judgment are avoided. Recorded ratings, prompts and review status support
inspection of the process without representing it as a human study.

Professional reporting also requires retaining unfavorable results. The matched control and
generic topics are therefore reported alongside positive coherence changes. Archived outputs
are not silently overwritten when data-quality concerns arise; changes to inclusion or training
would require a new versioned experiment and a renewed comparison.

\section{Discussion and Conclusions}
\label{sec:t1-conclusion}
The comparison illustrates a distinction between improving a metric and improving a useful
regional summary. Fine hyperparameter search improves several LDA diagnostics relative to the
grid baseline, but the observed GA does not outperform the matched random-search run. Neither
result establishes a universal ordering of optimisers. BERTopic provides more recognisable
domain labels under the automated rubric, at the cost of incomplete document assignment and
a coherence vocabulary that only partly overlaps the common reference dictionary.

For literature exploration, BERTopic's specific clusters are useful starting points for reading.
For summaries requiring a distribution for every record, LDA offers complete representation,
but its mixed topics should retain cautious labels. The practical recommendation is to inspect
topic words and representative documents alongside metrics and coverage, rather than allowing
the highest score to determine the interpretation automatically.

The limitations are substantive. The frozen corpus contains a known residual relevance error,
has sparse early coverage and does not exhaust local sources. Stability tests concern word lists
on the same corpus, automated ratings are not human validation, and only one complete run of
each search strategy is compared. Further work should adjudicate ambiguous records, rerun all
models on the revised corpus, repeat matched searches, examine sensitivity to fitness weights
and topic count, and obtain human judgments using the existing blinded instrument.

Within these limits, the study provides a reproducible comparison and evidence of increased
education-related prevalence in the later windows. Its main methodological lesson is that
coherence, interpretability, stability and coverage answer different questions. Making their
definitions and disagreements visible produces a more credible account of regional research
than selecting one model and treating its labels as established facts.

\balance
\section*{Individual Contributions}
Bibek Paudyal (250288) coordinated corpus retrieval, preprocessing and the topic-labelling
application. Sajan Mahat (250289) coordinated the LDA baseline, genetic search and matched
random-search experiments. Siddhartha Bhatta (250620) coordinated BERTopic, model evaluation
and temporal analysis. All three members contributed equally to report preparation and
review of the implementation, results and interpretation, with shared responsibility for
the final submission.

\begin{thebibliography}{99}
\bibitem{blei2003}
D.~M. Blei, A.~Y. Ng, and M.~I. Jordan, ``Latent Dirichlet allocation,'' \emph{Journal of Machine
Learning Research}, vol.~3, pp.~993--1022, 2003.

\bibitem{griffiths2004}
T.~L. Griffiths and M.~Steyvers, ``Finding scientific topics,'' \emph{Proc. National Academy of
Sciences}, vol.~101, suppl.~1, pp.~5228--5235, 2004.

\bibitem{yu2023}
D.~Yu and B.~Xiang, ``Discovering topics and trends in the field of artificial intelligence: Using
LDA topic modeling,'' \emph{Expert Systems with Applications}, 2023.

\bibitem{agrawal2018}
A.~Agrawal, W.~Fu, and T.~Menzies, ``What is wrong with topic modeling? (and how to fix it using
search-based software engineering),'' \emph{Information and Software Technology}, vol.~98,
pp.~74--88, 2018.

\bibitem{bergstra2012}
J.~Bergstra and Y.~Bengio, ``Random search for hyper-parameter optimization,'' \emph{Journal of
Machine Learning Research}, vol.~13, pp.~281--305, 2012.

\bibitem{grootendorst2022}
M.~Grootendorst, ``BERTopic: Neural topic modeling with a class-based TF-IDF procedure,''
\emph{arXiv:2203.05794}, 2022.

\bibitem{reimers2019}
N.~Reimers and I.~Gurevych, ``Sentence-BERT: Sentence embeddings using Siamese BERT-networks,'' in
\emph{Proc. EMNLP-IJCNLP}, 2019, pp.~3982--3992.

\bibitem{mcinnes2018}
L.~McInnes, J.~Healy, and J.~Melville, ``UMAP: Uniform manifold approximation and projection for
dimension reduction,'' \emph{arXiv:1802.03426}, 2018.

\bibitem{campello2013}
R.~J.~G.~B. Campello, D.~Moulavi, and J.~Sander, ``Density-based clustering based on hierarchical
density estimates,'' in \emph{Advances in Knowledge Discovery and Data Mining (PAKDD)}, 2013,
pp.~160--172.

\bibitem{agri2026}
J.~J. Capilitan, A.~L. Balbin, and J.~B. Matias, ``Mapping the evolution of artificial intelligence
in agriculture: A large-scale BERTopic analysis of smart farming, automation, and precision
systems (2020--2025),'' \emph{AgriEngineering}, vol.~8, no.~8, art.~312, 2026,
doi: 10.3390/agriengineering8080312.

\bibitem{chang2009}
J.~Chang, J.~Boyd-Graber, S.~Gerrish, C.~Wang, and D.~Blei, ``Reading tea leaves: How humans
interpret topic models,'' in \emph{Advances in Neural Information Processing Systems~22}, 2009.

\bibitem{newman2010}
D.~Newman, J.~H. Lau, K.~Grieser, and T.~Baldwin, ``Automatic evaluation of topic coherence,'' in
\emph{Proc. NAACL-HLT}, 2010, pp.~100--108.

\bibitem{mimno2011}
D.~Mimno, H.~M. Wallach, E.~Talley, M.~Leenders, and A.~McCallum, ``Optimizing semantic coherence
in topic models,'' in \emph{Proc. EMNLP}, 2011, pp.~262--272.

\bibitem{roder2015}
M.~R\"{o}der, A.~Both, and A.~Hinneburg, ``Exploring the space of topic coherence measures,'' in
\emph{Proc. 8th ACM Int. Conf. Web Search and Data Mining (WSDM)}, 2015, pp.~399--408.

\bibitem{dieng2020}
A.~B. Dieng, F.~J.~R. Ruiz, and D.~M. Blei, ``Topic modeling in embedding spaces,''
\emph{Transactions of the Association for Computational Linguistics}, vol.~8, pp.~439--453, 2020.

\bibitem{priem2022}
J.~Priem, H.~Piwowar, and R.~Orr, ``OpenAlex: A fully-open index of scholarly works, authors,
venues, institutions, and concepts,'' \emph{arXiv:2205.01833}, 2022.

\bibitem{rehurek2010}
R.~\v{R}eh\r{u}\v{r}ek and P.~Sojka, ``Software framework for topic modelling with large corpora,''
in \emph{Proc. LREC Workshop on New Challenges for NLP Frameworks}, 2010, pp.~45--50.

\bibitem{hoffman2010}
M.~D. Hoffman, D.~M. Blei, and F.~Bach, ``Online learning for latent Dirichlet allocation,'' in
\emph{Advances in Neural Information Processing Systems~23}, 2010, pp.~856--864.

\bibitem{holland1975}
J.~H. Holland, \emph{Adaptation in Natural and Artificial Systems}. Ann Arbor, MI: University of
Michigan Press, 1975.

\bibitem{krippendorff2004}
K.~Krippendorff, \emph{Content Analysis: An Introduction to Its Methodology}, 2nd ed. Thousand
Oaks, CA: Sage, 2004.

\end{thebibliography}


\section*{Appendix: Code and Reproducibility}
The implementation and experiment configuration are available at
\begin{center}
\small\url{https://github.com/bpdyl/ai-research-topic-modeling}
\end{center}
\noindent Pinned runs, source hashes and figure statistics are recorded in
\texttt{paper/report\_evidence.json}. Notebook quick mode checks execution; full settings
reproduce the protocol. Automated assessments use saved judgments.
